\section{Introduction}
\label{sec:intro}

A common pipeline in spectroscopic imaging reduces each pixel's spectrum
with a small encoder and classifies the resulting feature map with a wide
spatial network \citep{berisha2019deep,oleary2026spatial}. Hyperspectral
remote sensing, frozen image encoders under temporal models, and multimodal
fusion share the structure: a compact module sits behind a bottleneck, a
high-capacity module sits next to the loss, and both are trained jointly with
one learning rate. Practitioners often pretrain the small module and freeze
it. Whether freezing is necessary, or merely habitual, is exactly what a
fair comparison has to decide, and part of this paper is about what makes
such a comparison fair.

We ask what the loss geometry says about how the two modules learn, and what
a production model does and does not show. We distinguish three ingredients
of selective encoder adaptation: the curvature of the two parameter blocks,
the excitation of the encoder's directions by the training residual, and
the displacement the optimizer actually applies. Every statement we prove is
conditional, and we say on what.

\paragraph{Contributions.}
\begin{itemize}
\item \textbf{Conditional theory.} Under an explicit head hypothesis the
      module curvature disparity grows linearly in the downstream width
      (Theorem~\ref{thm:hessian}); under a residual envelope the encoder's
      displacement by the fitting time is
      $\mathcal{O}(\varepsilon\log(1/\delta))$ (Theorem~\ref{thm:twoscale});
      the encoder's share of the loss decrease is at most $a/(a+\kappa)$ at
      every horizon (Corollary~\ref{cor:attribution}); an affine bridge turns
      kernel coercivity on the visited residual into joint-versus-frozen
      closeness (Theorem~\ref{thm:bridge}), with a scoped nonlinear
      existence version (Section~\ref{sec:dynamics}).
\item \textbf{A solvable instance.} A serial cross-entropy model with
      replicated readouts has disparity exactly $M$, fitting rate $(M+1)/4$,
      a spectral update suppressed as $1/(1+Mv_0^2)$, and a failure under
      contextual reversal with probability $\Phi(m/2\sigma)$
      (Theorem~\ref{thm:instance}, Figure~\ref{fig:serial}); the displacement
      bound is within a factor of six of the truth.
\item \textbf{Diagnostic limits.} On a production spectral--spatial model the
      head hypothesis is structurally unmet, rank-one inputs invert the
      scalar disparity, the informative quantity is directional anisotropy,
      the observed width trend of downstream curvature is carried by
      batch-normalization gains, and on our own shallow instance the residual
      leaves the coercive subspace early (Section~\ref{sec:diagnostics},
      Figure~\ref{fig:diagnostics}).
\item \textbf{Matched protocol.} A freezing advantage observed under an
      earlier protocol is absent once clipping scope, normalization-parameter
      treatment and checkpoint policy are equalized
      (Section~\ref{sec:matched}, Figure~\ref{fig:matched},
      Table~\ref{tab:matched}); the original contrast cannot identify an
      architectural shortcut.
\end{itemize}

\paragraph{What the measurements teach.}
The theory is a set of sufficient conditions, and the production model
violates or evades most of them in instructive ways. Its classifier head has a
fixed input dimension, so the width-linear floor has no premise. Its inputs
are effectively rank one, so the top-eigenvalue ratio the theory bounds is
uninformative and the directional picture must replace it. The width growth
of its downstream curvature is produced by normalization gains acting on
width-diluted activations, which a function-preserving rescaling can move
freely. On the shallow instance where the geometry is proved, the residual
leaves the coercive part of the downstream kernel within a hundred steps, so
curvature at initialization predicts only an early phase. And in training,
three protocol asymmetries were enough to produce an apparent freezing
effect that a matched protocol does not reproduce. The paper's practical
content is therefore a set of measurements one must make before reading a
shortcut into a curvature ratio or a training curve, not a recipe.
